\begin{mistake}{2026-05-28}{08}

\begin{problemcontent}
设 \(f(x)\) 是连续函数，\(F(x)\) 是 \(f(x)\) 的原函数，则（ ）。

A. 当 \(f(x)\) 为奇函数时，\(F(x)\) 必为偶函数。

B. 当 \(f(x)\) 为偶函数时，\(F(x)\) 必为奇函数。

C. 当 \(f(x)\) 为周期函数时，\(F(x)\) 必为周期函数。

D. 当 \(f(x)\) 为单调函数时，\(F(x)\) 必为单调函数。
\end{problemcontent}

\begin{solutioncontent}
正确答案：A。

由原函数定义，\(F'(x)=f(x)\)。

A 项：若 \(f(x)\) 为奇函数，则 \(f(-x)=-f(x)\)。考察 \(F(-x)\) 与 \(F(x)\) 的关系：
\[
\frac{d}{dx}F(-x)=-F'(-x)=-f(-x)=f(x)=F'(x).
\]
因此 \(F(-x)-F(x)\) 为常数。令 \(x=0\)，得
\[
F(0)-F(0)=0,
\]
所以该常数为 \(0\)，即
\[
F(-x)=F(x).
\]
故 \(F(x)\) 必为偶函数，A 正确。

B 项：若 \(f(x)\) 为偶函数，\(\int_0^x f(t)\,dt\) 是奇函数，但一般原函数可加任意常数。取 \(f(x)=1\)，则 \(f(x)\) 为偶函数，原函数 \(F(x)=x+1\) 不是奇函数，故 B 错误。

C 项：取 \(f(x)=1\)，它是周期函数，但原函数 \(F(x)=x+C\) 不是周期函数，故 C 错误。

D 项：取 \(f(x)=x\)，它在 \(\mathbb{R}\) 上单调递增，但原函数
\[
F(x)=\frac{x^2}{2}+C
\]
在整个 \(\mathbb{R}\) 上不是单调函数，故 D 错误。

综上，正确选项为
\[
\boxed{A}
\]
\end{solutioncontent}

\mistakeknowledge{
\begin{itemize}
  \item \textbf{核心公式/定理}：若 \(F'(x)=f(x)\)，且 \(f\) 为奇函数，则 \(\frac{d}{dx}F(-x)=-f(-x)=f(x)=F'(x)\)，从而 \(F(-x)=F(x)\)。
  \item \textbf{易错点}：\(f\) 为偶函数时，\(\int_0^x f(t)\,dt\) 是奇函数，但一般原函数可加常数 \(C\)，不一定仍为奇函数。
  \item \textbf{关联知识}：周期函数的原函数不一定周期；原函数是否单调由 \(F'(x)=f(x)\) 的符号决定，而不是由 \(f(x)\) 是否单调决定。
\end{itemize}}

\end{mistake}
